% Preamble
\documentclass[twocolumn]{article}

% Packages
\usepackage[T1]{fontenc}
\usepackage[utf8]{inputenc}
\usepackage{indentfirst}
\usepackage{amsmath}
\usepackage[margin=1in]{geometry}
\usepackage[font=small,labelfont=bf]{caption}
\usepackage{hyperref}
\usepackage{fancyhdr}
\usepackage{booktabs}
\usepackage{graphicx}
\usepackage{listings}
\usepackage{xcolor}

\lstset{
    basicstyle=\ttfamily\footnotesize,
    breaklines=true,
    columns=fullflexible,
    keepspaces=true,
    language=C++,
    showstringspaces=false,
    frame=single,
    xleftmargin=0.5em,
}

\pagestyle{fancy}

\fancyhead{}
\lhead{
    \includegraphics[scale=0.13]{images/logo.png}
}
\rhead{
    Distributed and Parallel Computing \\
    2019/2020
    \footnotesize
}

\fancyfoot[C]{\thepage}

\date{}

\title{
    \includegraphics[scale=0.2]{images/logo.png}\\
    \vskip 0.5em
    \Huge Distributed and Parallel Computing\\
    \LARGE MPI Report --- 2D Grid\\
}

\author{
    Pedro Rio\\
    \textit{MSc Data Science and Engineering}\\
    \textit{Instituto Superior Técnico}\\
    \textit{97241}
    \and
    Sebastião Maia Serqueira\\
    \textit{MSc Data Science and Engineering}\\
    \textit{Instituto Superior Técnico}\\
    \textit{97108}\\
    \and
    Eduardo Machado\\
    \textit{MSc Data Science and Engineering}\\
    \textit{Instituto Superior Técnico}\\
    \textit{97133}
}

% Document
\begin{document}
    \maketitle


    \section{Program Overview}\label{sec:program-overview}
    The problem is a matrix-factorisation recommender: a sparse user--item matrix $A$ is
    approximated by $A \approx L \cdot R$, with $L$ of size
    $\textit{users} \times \textit{features}$ and $R$ of size
    $\textit{features} \times \textit{items}$. The factors are fitted by gradient descent
    over the observed entries, and the final dense product $B = L \cdot R$ yields one
    recommendation per user. The four stages are \emph{ReadInput}, \emph{InitialLR},
    \emph{UpdateLR} (repeated for the requested number of iterations), and
    \emph{FilterFinalMatrix}.


    \section{Why a 2D Grid}\label{sec:why-grid}
    The simplest MPI design is \emph{replicate-all}: every process keeps full copies of
    $L$ and $R$, splits the non-zero list into blocks, computes a partial gradient, and
    sums it with a global \texttt{MPI\_Allreduce}. It is correct and short, but it does not
    scale: the reduced gradient has $O(\textit{features} \cdot (\textit{users} +
    \textit{items}))$ entries and that volume \emph{does not shrink} as processes are
    added, while compute per process falls as $1/P$. Communication therefore dominates at
    scale, and on \emph{wide} instances (a large $R$) it \emph{anti-scales} --- adding
    processes makes it slower. Memory does not scale either: a single node must hold the
    whole problem.

    Both limits have the same cause --- replication --- so we distribute the factors over a
    2D process grid.


    \section{The 2D-Grid Version}\label{sec:grid}

    \subsection{Process mesh}\label{subsec:mesh}
    \texttt{MPI\_Dims\_create} factors $P$ into a balanced $P_r \times P_c$ grid;
    \texttt{MPI\_Cart\_create} builds the mesh (with reorder disabled, so a rank's grid id
    equals its world rank) and \texttt{MPI\_Cart\_sub} derives a \emph{row} and a
    \emph{column} sub-communicator for each rank.

    \subsection{Foster methodology}\label{subsec:foster}
    \textbf{Partitioning.} $L$ is split by \emph{user-rows} across the grid rows and $R$ by
    \emph{item-columns} across the grid columns. A non-zero $(u,i)$ is owned by rank
    $(\textit{userBlock}(u), \textit{itemBlock}(i))$, so each rank holds exactly the
    $L$-rows and $R$-columns its non-zeros touch --- memory per rank is
    $O((\textit{users} \cdot \textit{features} + \textit{features} \cdot
    \textit{items})/\sqrt{P})$.

    \textbf{Communication.} Each rank computes a partial gradient from its cell's
    non-zeros (all data local), then $dL$ is summed across the \emph{row} communicator and
    $dR$ down the \emph{column} communicator. These are reductions of $\sqrt{P}$-smaller
    slices over $\sqrt{P}$-sized sub-communicators, not a global reduce of the whole
    gradient.

    \textbf{Agglomeration and Mapping.} The non-zeros are agglomerated into one contiguous
    block per grid cell and routed there once with a bucketed \texttt{MPI\_Scatterv}; one
    cell maps to one process.

    \subsection{Deterministic, communication-free init}\label{subsec:init}
    \texttt{initialLR} is seeded with a fixed value, but a partitioned factor cannot simply
    seek into the random stream (an $R$-column slice is strided). The trick is
    \emph{generate-and-discard}: every rank runs the same \texttt{srandom(1)} sequence in
    the serial fill order but stores only the values that fall in its slice. The slices
    hold exactly the serial values, with no broadcast of the factors.

    \subsection{Distributed final stage}\label{subsec:final}
    $B = L \cdot R$, the mask of known entries, and the per-user arg-max are all computed
    on the cells: each rank forms its $B$-block, takes a local arg-max over its item
    columns, and the results are combined across the row with \texttt{MPI\_MAXLOC}
    (smallest item index on ties, matching the serial scan) and gathered to the root in
    user order. There is no serial-root tail.


    \section{A Note on the Hybrid Variant}\label{sec:hybrid}
    This report covers the \textbf{pure MPI} deliverable: all parallelism comes from MPI
    processes, each rank runs its share sequentially, and the binary links no OpenMP
    runtime. A second deliverable, \texttt{mpi-hybrid/}, keeps this same 2D-grid
    decomposition and adds OpenMP threads inside each rank; it has its own report. Every
    figure below was measured with the pure variant.


    \section{Evaluation}\label{sec:evaluation}
    Table~\ref{tab:grid-vs-rep} compares wall-clock time of the 2D-grid version against
    the naive replicate-all design it replaced.

    \textbf{Provenance.} The \textbf{grid} rows are the committed measurements in
    \texttt{compare/data/comparison.mpi.csv}, produced by
    \texttt{compare/scripts/benchmark.py} on an Apple Silicon laptop (macOS 15.6, 10
    cores, 8 of them performance cores), three repetitions with the minimum kept, every
    run verified against the instance's \texttt{.out} file. They are reproducible from
    this repository by one command. The \textbf{replicate} rows were measured on the same
    machine before the rewrite and \emph{cannot} be regenerated here: that implementation
    was replaced and no longer exists in the tree. They are retained because the
    comparison is the reason the 2D grid exists, and labelled so no one mistakes them for
    current figures.

    \begin{table}[!h]
        \caption{\label{tab:grid-vs-rep}Replicate-all vs.\ 2D-grid (time, s). Grid rows:
        measured now, this laptop. Replicate rows: pre-rewrite, same machine, not
        regenerable.}
        \centering
        \resizebox{\linewidth}{!}{
            \begin{tabular}[t]{llrrrr}
                \toprule
                instance & design & 1 & 2 & 4 & 8\\
                \midrule
                inst200 (wide) & replicate\textsuperscript{*} & 7.62 & 5.26 & 5.36 & 10.64\\
                               & \textbf{grid} & 6.07 & 3.94 & 2.25 & 2.94\\
                \addlinespace
                inst500        & replicate\textsuperscript{*} & 8.50 & 4.84 & 3.28 & 8.41\\
                               & \textbf{grid} & 8.48 & 4.76 & 2.66 & 3.30\\
                \addlinespace
                inst1000       & replicate\textsuperscript{*} & 4.15 & 2.70 & 1.89 & 3.03\\
                               & \textbf{grid} & 3.96 & 2.22 & 1.23 & 0.99\\
                \bottomrule
            \end{tabular}}

        \footnotesize\textsuperscript{*}pre-rewrite implementation, not regenerable from
        this repository.
    \end{table}

    The grid version is faster at every process count on every instance, and the gap
    widens with $P$: at 8 processes it is $3.6\times$ faster on the wide instance and
    $3.1\times$ on inst1000. Replicate-all \emph{regresses} sharply past 4 processes ---
    $10.64$\,s on the wide instance, worse than its own single-process time --- because
    the dense \texttt{Allreduce} of the whole gradient becomes a communication ceiling.

    The grid version is not immune to the 8-process point either: it regresses on the wide
    instance ($2.25 \to 2.94$\,s) and on inst500 ($2.66 \to 3.30$\,s), while still
    improving on inst1000 ($1.23 \to 0.99$\,s). Two effects overlap there, and this
    measurement cannot separate them: the machine has 10 cores but only 8 performance
    cores, so an 8-process run competes with the operating system and the timing harness
    for the last of them; and the per-process work at that point is small enough that
    communication is a larger share of it. The honest claim is therefore that the grid
    design \emph{degrades far more gracefully}, not that it scales indefinitely --- a
    $2.94$\,s regression against $10.64$\,s is a difference in kind, but it is still a
    regression.

    \section{Implications for Parallel and Distributed Computing}\label{sec:pdc}

    \subsection{Distribute the data, not just the work}\label{subsec:distribute}
    Replicate-all distributes the \emph{work} but not the \emph{data}, so it hits both a
    memory wall and a communication ceiling. The 2D grid distributes the data: the factor
    slices and the gradient fall as $1/\sqrt{P}$, and the dense $B$ built in the final
    stage --- which actually dominates the footprint on every instance here, by roughly a
    factor of five --- falls as $1/P$. Replicate-all's per-rank memory does not fall at
    all. The general lesson is that a scalable design must shrink each process's
    footprint, not only its share of the compute; and that the term worth quoting is the
    one that dominates, not the one that reads best.

    \subsection{Match the collective to the decomposition}\label{subsec:collective}
    The win is not a faster \texttt{Allreduce}; it is a \emph{smaller} one, over a
    \emph{smaller} communicator. Because a user's updates are spread only across its grid
    row and an item's only down its grid column, the global reduction decomposes into
    independent row and column reductions. Choosing the communicator to match the data
    layout is what turns flat communication into communication that scales.

    \subsection{Remove the serial tail}\label{subsec:tail}
    Computing $B$ and the arg-max on the root would be an
    $O(\textit{users} \cdot \textit{items} \cdot \textit{features})$ sequential tail that,
    by Amdahl's law, caps the speedup. Distributing it over the same grid keeps the whole
    pipeline parallel.

    \subsection{Reproducibility under reordering}\label{subsec:repro}
    The partitioned reductions sum the gradient in a different order than the serial code,
    so floating-point results differ in the last bits --- but the arg-max is robust to
    that, so the recommendations match the serial output byte for byte at every process
    count and grid shape. A deterministic serial reference is what makes this verifiable,
    and the test suite asserts it rather than assuming it: 32 cases covering every
    instance at 1, 2 and 4 processes, non-square and prime grids, and blocks left empty
    when processes outnumber users or items.

\end{document}
